% TEMPLATE-MILC-v3.tex - plantilla maestra UNICA de las guias MILC v3.
%
% La usan las tres familias de guias, cada una con su driver:
%   · Grados 6-11  -> scripts/build-guias-g11.py
%   · Bebras       -> scripts/build-guias-bebras.py
%   · Semillero    -> scripts/build-guias-semillero.py
%
% Antes habia tres copias casi identicas (~400 lineas iguales, 6 distintas):
% cada mejora habia que aplicarla tres veces, y en la practica no se hacia
% (el motor de recursos vivia solo en dos de ellas). Lo que varia entre
% familias esta parametrizado en 6 placeholders que cada driver llena
% (se nombran aqui SIN su sintaxis real: el builder reemplaza por texto plano,
% tambien dentro de los comentarios, y un valor multilinea rompe el preambulo):
%
%   HEADER_IZQ        encabezado izquierdo de pagina
%   HEADER_DER        encabezado derecho de pagina
%   PORTADA_ETIQUETA  rotulo sobre el numero grande (GRADO/MOMENTO/MODULO)
%   PORTADA_SERIE     linea de serie de la portada
%   PORTADA_ID        identificador de la guia en la portada
%   PORTADA_CIRCULO   el circulito de grado (°), o vacio si no aplica
%
% El resto de claves las documenta content/guias/_SCHEMA.md. El script valida
% que no queden placeholders sin reemplazar antes de escribir el archivo final.

\documentclass[11pt,letterpaper]{article}
\usepackage[margin=2.0cm]{geometry}
\usepackage[table]{xcolor}
\usepackage{fontspec}
\usepackage[spanish]{babel}
\usepackage{tikz}
% Librerías TikZ para los diagramas de las guías (bloque `recursos:`):
%  · babel  → babel en español vuelve activos caracteres como «>», y TikZ los usa
%             en sus opciones (>=stealth, ->). Sin esta librería, cualquier
%             diagrama con flechas revienta con «\language@active@arg> has an
%             extra }». Debe ir ANTES de que se dibuje nada.
%  · positioning        → sintaxis «below left=2pt and 3pt».
%  · decorations.pathreplacing → llaves ({brace}) para señalar rangos.
\usetikzlibrary{babel,positioning,decorations.pathreplacing}
\usepackage{graphicx}
% Circuitos: el trabajo de electrónica se hace hoy con micro:bit y sus sensores
% integrados (MakeCode, sin cableado), así que no se necesitan esquemáticos.
% Si algún día entran montajes con protoboard, basta descomentar esta línea:
% los diagramas .tex/.tikz declarados en `recursos:` ya se incrustan solos.
% \usepackage{circuitikz}
\usepackage[most]{tcolorbox}
\usepackage{tabularx}
\usepackage{array}
\usepackage{fancyhdr}
\usepackage{titlesec}
\usepackage[document]{ragged2e}  % bandera a la derecha en todo el cuerpo
\usepackage{enumitem}
\usepackage{microtype}

% Helvetica Neue no trae la flecha U+2192, asi que cada «→» escrito literalmente
% en el YAML se compilaba a NADA, en silencio (462 flechas en 61 guias: rutas del
% tipo «paleta → tipografia → lockup» salian rotas). Mapeamos el caracter a su
% version matematica, que si tiene glifo. Asi el autor puede seguir escribiendo
% «→» comodamente en el YAML.
\usepackage{newunicodechar}
\newunicodechar{→}{\ensuremath{\rightarrow}}
% Mismo problema con los simbolos de marca y vinetas: el «✓» de «una palomita ✓»
% salia como cuadrito vacio en el PDF de todas las guias que lo usaban.
\usepackage{pifont}
\newunicodechar{✓}{\ding{51}}
\newunicodechar{✗}{\ding{55}}
\newunicodechar{✔}{\ding{52}}
\newunicodechar{•}{\textbullet}
\newunicodechar{≥}{\ensuremath{\geq}}
\newunicodechar{≤}{\ensuremath{\leq}}

\setmainfont{Helvetica Neue}
\setsansfont{Helvetica Neue}
\setlength{\parindent}{0pt}
\setlength{\parskip}{6pt}
% Sin particion de palabras: en 6.º-7.º una palabra cortada por guion al final
% de linea («Comuni-cacion») frena la lectura mucho mas que una linea corta.
\hyphenpenalty=10000
\exhyphenpenalty=10000
\pretolerance=2000
\tolerance=3000
\setlength{\emergencystretch}{3em}
\linespread{1.10}
\renewcommand{\arraystretch}{1.25}
\setlist{leftmargin=5mm,itemsep=1mm,topsep=1mm}
\newcolumntype{Y}{>{\RaggedRight\arraybackslash}X}

\definecolor{milcMagenta}{HTML}{E6007E}
\definecolor{milcVino}{HTML}{5A0038}
\definecolor{milcTurquesa}{HTML}{0093A5}
\definecolor{milcVerde}{HTML}{58A923}
\definecolor{milcMostaza}{HTML}{E5B400}
\definecolor{milcOcre}{HTML}{B86600}
\definecolor{milcNegro}{HTML}{241816}
\definecolor{milcGris}{HTML}{F7F7F4}
\definecolor{milcLinea}{HTML}{E8E2DD}
\definecolor{milcGradePrimary}{HTML}{84CC16}
\definecolor{milcGradeDark}{HTML}{4D7C0F}
\definecolor{milcGradeSoft}{HTML}{ECFCCB}
\definecolor{milcGradeAccent}{HTML}{CFE7FF}
\definecolor{milcGradeText}{HTML}{FFFFFF}
\color{milcNegro}

\pagestyle{fancy}
\fancyhf{}
\lhead{\small Tecnología e Informática · Grado 7}
\rhead{\small Guía 25 · MILC v3}
\cfoot{\small\thepage}
\renewcommand{\headrulewidth}{0pt}

\newtcolorbox{titlebox}[1]{
  enhanced,colback=#1,colframe=#1,arc=7mm,boxrule=0pt,
  left=7mm,right=7mm,top=3mm,bottom=3mm,
  before skip=8pt,after skip=8pt,
  fontupper=\sffamily\bfseries\Large\color{white}
}

\newtcolorbox{softbox}[3]{
  enhanced,colback=#2,colframe=#1,borderline west={4pt}{0pt}{#1},
  arc=2mm,boxrule=.4pt,left=4mm,right=4mm,top=3mm,bottom=3mm,
  before skip=6pt,after skip=8pt,title={#3},coltitle=white,
  colbacktitle=#1,fonttitle=\sffamily\bfseries,fontupper=\RaggedRight,
  attach boxed title to top left={xshift=3mm,yshift=-2mm},
  boxed title style={arc=2mm, boxrule=0pt}
}

\newtcolorbox{drawbox}[2]{
  enhanced,colback=white,colframe=#1,arc=4mm,boxrule=1pt,
  left=5mm,right=5mm,top=4mm,bottom=4mm,before skip=8pt,after skip=8pt,
  title={#2},coltitle=white,colbacktitle=#1,fonttitle=\sffamily\bfseries,
  fontupper=\RaggedRight,
  attach boxed title to top left={xshift=4mm,yshift=-2mm},
  boxed title style={arc=3mm, boxrule=0pt}
}

\newtcolorbox{infoband}[1]{
  enhanced,colback=#1!8,colframe=#1!50,arc=2mm,boxrule=.5pt,
  left=4mm,right=4mm,top=2mm,bottom=2mm,before skip=4pt,after skip=4pt,
  fontupper=\small\RaggedRight
}

% Puente narrativo entre secciones (contrato editorial Conéctate).
% Da continuidad: "Acabas de X. Ahora vas a Y porque Z."
% Visualmente: banda izquierda en color del grado, texto en itálica pequeña.
\newtcolorbox{puentebox}{
  enhanced,colback=milcGradeSoft,colframe=milcGradeDark,
  borderline west={3pt}{0pt}{milcGradeDark},
  arc=0mm,boxrule=0pt,left=5mm,right=4mm,top=2.5mm,bottom=2.5mm,
  before skip=4pt,after skip=8pt,
  fontupper=\itshape\small\color{milcNegro}\RaggedRight
}
% La flecha va en modo matematico a proposito: Helvetica Neue NO tiene el glifo
% U+2192, asi que \textrightarrow se compilaba a NADA («Missing character: There
% is no → in font Helvetica Neue») y el puente salia sin flecha. En modo
% matematico la flecha viene de la fuente matematica y si se dibuja.
\newcommand{\puente}[1]{%
  \begin{puentebox}%
    \textcolor{milcGradeDark}{$\rightarrow$}\hspace{2mm}#1%
  \end{puentebox}%
}

\newcommand{\linead}[1]{\noindent\color{milcLinea}\rule{#1}{0.5pt}\color{milcNegro}}
\newcommand{\respuesta}[1]{\par\vspace{2mm}\foreach \n in {1,...,#1}{\noindent\color{milcLinea}\rule{\linewidth}{0.5pt}\par\vspace{5mm}}\color{milcNegro}}
\newcommand{\checkbox}{\textcolor{milcGradeDark}{\fbox{\phantom{x}}}\hspace{5pt}}

\newcommand{\phasecard}[4]{
\begin{tcolorbox}[enhanced,colback=#3,colframe=#2,arc=3mm,boxrule=.5pt,height=3.05cm,valign=center,halign=center,left=1.5mm,right=1.5mm,top=2mm,bottom=2mm]
{\sffamily\bfseries\fontsize{22}{24}\selectfont\textcolor{#2}{#1}}\par
{\sffamily\bfseries\small #4}
\end{tcolorbox}
}

% ── Piezas del rediseno 2026 (legibilidad 6.º-11.º) ───────────────────
% Banda de estacion: reemplaza el titlebox macizo. Titulo a la izquierda,
% dato practico (tiempo/modalidad) a la derecha, en una sola linea.
\newcommand{\milcestacion}[2]{%
\begin{tcolorbox}[enhanced,colback=#1,colframe=#1,arc=2mm,boxrule=0pt,
  left=5mm,right=5mm,top=2.6mm,bottom=2.6mm,before skip=11pt,after skip=7pt]
{\sffamily\bfseries\fontsize{15}{18}\selectfont\color{white}#2}
\end{tcolorbox}}

% Tarjeta de paso numerada: sustituye las tablas «Pilar / Aplicacion», que
% eran formato de consulta docente, no de estudiante. El numero va en un
% circulo sobre el borde para que la secuencia se lea de un vistazo.
\newcommand{\milcpaso}[3]{%
\begin{tcolorbox}[enhanced,colback=white,colframe=milcLinea,arc=1.5mm,boxrule=.9pt,
  left=14mm,right=4mm,top=3mm,bottom=3mm,before skip=4pt,after skip=4pt,
  fontupper=\RaggedRight,
  overlay={\node[anchor=north west,circle,fill=#2,text=white,inner sep=0pt,
    minimum size=7.4mm,font=\sffamily\bfseries\small]
    at ([xshift=3.6mm,yshift=-3.2mm]frame.north west) {#1};}]
#3
\end{tcolorbox}}

% Casilla marcable de tamano util con lapiz (la anterior era \fbox{\phantom{x}},
% de alto variable segun el texto que la seguia).
\newcommand{\milcmarca}{\framebox[3.8mm]{\rule{0pt}{3.4mm}}}

% Fila marcable de lista suelta (checks de la estacion 1).
\newcommand{\milccheckrow}[1]{%
\par\vspace{1.8mm}\noindent
\makebox[6.4mm][l]{\raisebox{-.55mm}{\textcolor{milcNegro}{\milcmarca}}}%
\parbox[t]{\dimexpr\linewidth-6.4mm\relax}{\RaggedRight #1}\par}

% Celda marcable dentro de la rubrica (tabla de autoevaluacion). Misma
% sangria francesa, pero pensada para vivir dentro de una columna X.
\newcommand{\milcopcion}[1]{%
\makebox[5.2mm][l]{\raisebox{-.55mm}{\textcolor{milcNegro}{\milcmarca}}}%
\parbox[t]{\dimexpr\linewidth-5.2mm\relax}{\RaggedRight #1}}

% ── Recursos visuales de la guía ──────────────────────────────────────
% Declarados en el bloque `recursos:` del YAML y emitidos por el builder.
% \guiaFigura   → imagen rasterizada o PDF (foto, carta celeste, montaje).
% \guiaDiagrama → fuente TikZ/circuitikz (.tex) incrustada como vector:
%                 es la vía para circuitos y esquemáticos editables.
% #1 = ruta relativa al .tex · #2 = pie (puede ir vacío)
\newcommand{\guiaFigura}[2]{%
\begin{center}
\includegraphics[width=0.88\linewidth,height=0.38\textheight,keepaspectratio]{#1}\\[1.5mm]
{\footnotesize\itshape\textcolor{milcNegro!75}{#2}}
\end{center}
}

\newcommand{\guiaDiagrama}[2]{%
\begin{center}
\input{#1}\\[1.5mm]
{\footnotesize\itshape\textcolor{milcNegro!75}{#2}}
\end{center}
}

\begin{document}
\thispagestyle{empty}

% ── Portada ───────────────────────────────────────────────────────────
% Antes: un rectangulo de color a pagina completa. Se veia bien en pantalla,
% pero en la fotocopiadora del colegio se comia el toner y salia gris sucio.
% Ahora la banda ocupa el tercio superior y el resto va en blanco.
\begin{tikzpicture}[remember picture,overlay]
  \path[fill=milcGradePrimary]
    (current page.north west) rectangle ([yshift=-10.6cm]current page.north east);

  \node[anchor=north west,text=milcGradeText] at ([xshift=2.0cm,yshift=-1.55cm]current page.north west)
    {\sffamily\bfseries\fontsize{12}{14}\selectfont GRADO};
  \node[anchor=north west,text=milcGradeText,opacity=.85] at ([xshift=2.0cm,yshift=-2.35cm]current page.north west)
    {\sffamily\bfseries\fontsize{8.5}{10}\selectfont SERIE GUÍAS · TECNOLOGÍA E INFORMÁTICA};
  \node[anchor=north east,text=milcGradeText,opacity=.85,align=right] at ([xshift=-2.0cm,yshift=-1.55cm]current page.north east)
    {\sffamily\bfseries\fontsize{9}{12}\selectfont I.E. Sor María Juliana\\Cartago, Valle del Cauca};

  \node[anchor=west,text=milcGradeText] (gradonum) at ([xshift=1.85cm,yshift=-7.1cm]current page.north west)
    {\sffamily\bfseries\fontsize{112}{112}\selectfont 7};
  % El circulito de grado (°) solo aplica a las guias de grado («11°»); en
  % Bebras y Semillero el numero grande es un momento/modulo, no un grado.
  % Se posiciona relativo al nodo `gradonum`, no en coordenadas absolutas.
  \draw[milcGradeText,line width=1.6pt]
    ([xshift=2.0mm,yshift=-4.5mm]gradonum.north east) circle (.15cm);

  \node[anchor=west,text=milcGradeText,text width=11.4cm,align=left] at ([xshift=1.5cm,yshift=0.5cm]gradonum.east)
    {
     {\sffamily\bfseries\fontsize{9}{11}\selectfont\MakeUppercase{Período 3 · Inteligencia Artificial}}\\[3mm]
     {\sffamily\bfseries\fontsize{21}{25}\selectfont\setlength{\baselineskip}{27pt}Modelos de\\[4pt]lenguaje\\[4pt](LLMs) ---\\[4pt]ChatGPT,\\[4pt]Claude, Gemini}\\[3.5mm]
     {\sffamily\bfseries\fontsize{15}{18}\selectfont Guía 25-7-TIC}
    };
\end{tikzpicture}

\vspace*{9.4cm}
\vfill

{\sffamily\bfseries\fontsize{8}{10}\selectfont\textcolor{milcGradeDark}{LO QUE TE LLEVAS HOY}}

\begin{tcolorbox}[enhanced,colback=white,colframe=milcNegro,arc=2mm,boxrule=1.6pt,
  left=5mm,right=5mm,top=4mm,bottom=4mm,before skip=3pt,after skip=12pt,
  fontupper=\RaggedRight\fontsize{11}{15.5}\selectfont]
\textbf{3 conversaciones reales con un LLM} gratuito (puedes usar ChatGPT, Claude, Gemini o Copilot), guardadas en capturas o transcritas en el cuaderno. \textbf{Tema sugerido para las 3:} \emph{(1) ``Explícame algo que me cueste entender de matemáticas''}, \emph{(2) ``Recomiéndame 3 libros de un género que me guste''}, \emph{(3) ``Ayúdame a redactar un correo formal al profe''}. Más \textbf{tabla comparativa} (en cuaderno) de 3 LLMs principales (ChatGPT vs Claude vs Gemini) + \textbf{análisis} de las 5 cosas que SÍ pueden y las 5 que NO pueden hacer.
\end{tcolorbox}

\begin{tcolorbox}[enhanced,colback=milcGris,colframe=milcGris,arc=2mm,boxrule=0pt,
  left=5mm,right=5mm,top=3.5mm,bottom=3.5mm,after skip=12pt]
{\sffamily\bfseries\fontsize{8}{10}\selectfont\textcolor{milcNegro!55}{ESTA GUÍA ES DE}}\\[3mm]
\begin{tabularx}{\linewidth}{X p{3.2cm} p{3.2cm}}
\linead{\linewidth} & \linead{3.0cm} & \linead{3.0cm}\\[-1mm]
{\fontsize{8.5}{10}\selectfont\textcolor{milcNegro!55}{Nombre completo}} &
{\fontsize{8.5}{10}\selectfont\textcolor{milcNegro!55}{Grupo}} &
{\fontsize{8.5}{10}\selectfont\textcolor{milcNegro!55}{Fecha}}\\
\end{tabularx}
\end{tcolorbox}

\vfill

\noindent\textcolor{milcLinea}{\rule{\linewidth}{1pt}}\\[2mm]
{\sffamily\bfseries\fontsize{9.5}{12}\selectfont Autor: Dr. Álvaro Cárdenas Orozco}\\[1mm]
{\sffamily\fontsize{8.5}{11}\selectfont\textcolor{milcNegro!55}{Metodología MILC · apertura ancestral · pensamiento computacional · triángulo Dussel–estoicismo–Floridi}}

\newpage

\section*{Apertura: Modelos de lenguaje (LLMs) --- ChatGPT, Claude, Gemini}

\begin{softbox}{milcGradeDark}{milcGradeSoft}{Saber ancestral}
\textbf{Cuando doña Mercedes la maestra rural de Cartago tenía que explicarle un tema difícil a un niño del campo, no le daba la respuesta de una vez.} \emph{Le hacía preguntas, le mostraba ejemplos, lo guiaba paso a paso}. Si el niño no entendía con la primera explicación, doña Mercedes lo intentaba de otra manera: con un dibujo, con una comparación con la vida del campo, con un cuento. \textbf{Era paciente y se adaptaba al estudiante}. Los LLMs (Large Language Models) intentan hacer algo parecido: \emph{conversan contigo, se adaptan a tus preguntas, te explican de varias formas}. Pero hay diferencias claves: \emph{doña Mercedes te conocía, sabía qué te interesaba, recordaba tu progreso. Los LLMs no te conocen}: cada conversación arranca de cero. Esa diferencia es lo que hace que la IA sea \emph{asistente} pero no \emph{maestra} en sentido completo.
\end{softbox}

\begin{softbox}{milcTurquesa}{milcGris}{Saber contemporáneo}
Los \textbf{LLMs} (Large Language Models, \emph{Modelos de Lenguaje Grandes}) son IAs entrenadas con \textbf{miles de millones de palabras} de libros, sitios web, foros, artículos. Aprenden a \emph{predecir la siguiente palabra} con tanta sofisticación que pueden mantener conversaciones, responder preguntas, escribir textos, traducir, resumir. \textbf{Los 3 LLMs más conocidos en 2026:} \textbf{(1) ChatGPT} (de OpenAI): el primero masivo (nov 2022). Versión gratuita: GPT-4o-mini. Pago: GPT-4o, o3. \textbf{(2) Claude} (de Anthropic): conocido por respuestas cuidadosas y análisis profundo. Versión gratuita disponible. \textbf{(3) Gemini} (de Google): integrado con servicios de Google. Versión gratuita en \texttt{gemini.google.com}. \textbf{Cosas que pueden hacer (las 5 más útiles para estudiantes):} explicar conceptos, ayudar a estructurar texto, traducir, resumir, generar ideas. \textbf{Cosas que NO pueden (5 limitaciones clave):} (1) no tienen \emph{información actualizada} más allá de su fecha de entrenamiento; (2) \emph{alucinan} (inventan); (3) no acceden a internet en tiempo real (excepto algunas versiones); (4) no tienen \emph{memoria} entre conversaciones; (5) no tienen \emph{contexto personal} sobre ti.
\end{softbox}

\begin{softbox}{milcMostaza}{milcGris}{Pregunta-puente}
\textit{Si le preguntas a ChatGPT \emph{``¿quién ganó las elecciones de mi colegio el año pasado?''}, ¿\textbf{podrá responder}? ¿Por qué sí o por qué no?}
\end{softbox}

\begin{softbox}{milcVerde}{milcGris}{Saber y saber hacer}
Al terminar la clase: (1) podrás \textbf{identificar} los 3 LLMs principales y sus diferencias; (2) sabrás \textbf{aplicar} cada uno a tareas reales; (3) podrás \textbf{evaluar} qué pueden y qué NO pueden hacer; (4) habrás \textbf{conversado} con un LLM real al menos 3 veces.
\end{softbox}



\small
\begin{tabularx}{\textwidth}{>{\bfseries}p{3.0cm}Y}
\rowcolor{milcGradePrimary!12}Autor & Dr. Álvaro Cárdenas Orozco\\
DBA & Reconoce los fundamentos, usos y límites éticos de la Inteligencia Artificial (MEN, T\&I 7°)\\
Referentes & Saber ancestral del consejero · Modelos de lenguaje · Ética algorítmica y prompting\\
Producto final & \textbf{3 conversaciones reales con un LLM} gratuito (puedes usar ChatGPT, Claude, Gemini o Copilot), guardadas en capturas o transcritas en el cuaderno. \textbf{Tema sugerido para las 3:} \emph{(1) ``Explícame algo que me cueste entender de matemáticas''}, \emph{(2) ``Recomiéndame 3 libros de un género que me guste''}, \emph{(3) ``Ayúdame a redactar un correo formal al profe''}. Más \textbf{tabla comparativa} (en cuaderno) de 3 LLMs principales (ChatGPT vs Claude vs Gemini) + \textbf{análisis} de las 5 cosas que SÍ pueden y las 5 que NO pueden hacer.\\
\end{tabularx}
\normalsize

\puente{Hoy haces 4 cosas: distingues los 3 LLMs principales, abres uno y conversas 3 veces, evalúas sus respuestas, analizas qué pueden y qué no.}

\milcestacion{milcGradeDark}{Ruta MILC: así vamos a aprender}
\begin{tabularx}{\textwidth}{>{\centering\arraybackslash}X>{\centering\arraybackslash}X>{\centering\arraybackslash}X>{\centering\arraybackslash}X}
\phasecard{1}{milcVerde}{milcGris}{Escucha\\[-1mm]\small Observo, escucho y pregunto.} &
\phasecard{2}{milcTurquesa}{milcGris}{Entender\\[-1mm]\small Sistematizo y ordeno.} &
\phasecard{3}{milcMagenta}{milcGris}{Hacer\\[-1mm]\small Construyo y aplico.} &
\phasecard{4}{milcVino}{milcGris}{Cierre\\[-1mm]\small Evalúo y reflexiono.}
\end{tabularx}


\puente{Empezamos planeando las 3 conversaciones.}

\milcestacion{milcVerde}{Estación 1 de 3 · Escucha}

\begin{softbox}{milcVerde}{milcGris}{Escena para escuchar}
\textbf{Actividad 1 · IDENTIFICA --- Planea tus 3 conversaciones} (10 min · individual). En el cuaderno, planea las 3 preguntas que vas a hacerle al LLM. \textbf{Conversación 1: Aprender algo difícil.} Escoge un concepto de matemáticas o ciencias que te cueste. Anota: \emph{``Voy a pedirle que me explique [concepto] de manera fácil''}. \textbf{Conversación 2: Recomendaciones.} Anota un género de libros que te guste (aventura, ciencia ficción, romance) y pide 3 títulos. \textbf{Conversación 3: Redactar algo.} Piensa en un correo formal que tengas que escribir (al profe, al coordinador). Anota qué quieres pedir y a quién.
\end{softbox}

{\sffamily\bfseries\fontsize{8}{10}\selectfont\textcolor{milcNegro!55}{MARCA CUANDO LO HAYAS HECHO}}
\milccheckrow{Pensé en un concepto que me cueste para Conv. 1.}
\milccheckrow{Identifiqué mi género favorito para Conv. 2.}
\milccheckrow{Definí qué correo formal necesito para Conv. 3.}

\begin{infoband}{milcVerde}


\textbf{Lo que va al cuaderno (Actividad 1 · IDENTIFICA).} Título: «Actividad 1 --- Plan de mis 3 conversaciones». \textbf{Formato:} 3 bloques con la pregunta planeada para cada conversación. \textbf{Extensión:} 6-9 líneas. \textbf{Sabes que terminaste cuando} tienes claro qué le vas a preguntar al LLM antes de abrirlo.
\end{infoband}

\puente{Después aprendes los 3 LLMs principales + sus 5 capacidades y 5 limitaciones.}

\milcestacion{milcTurquesa}{Estación 2 de 3 · Entender}

\begin{softbox}{milcTurquesa}{milcGris}{Lo que vas a entender}
Antes de abrir el LLM, conoce los 3 principales + sus 5 capacidades + sus 5 limitaciones. Eso te coloca en posición de evaluar las respuestas, no solo aceptarlas.
\end{softbox}

\milcpaso{1}{milcTurquesa}{\textbf{Los 3 LLMs principales en 2026.} \textbf{(1) ChatGPT (OpenAI).} El primer LLM masivo (nov 2022). Versión gratuita: GPT-4o-mini en \texttt{chat.openai.com}. Conocido por ser muy versátil. Tiene memoria entre conversaciones en cuenta paga. \textbf{(2) Claude (Anthropic).} Conocido por respuestas cuidadosas, largas y bien estructuradas. Versión gratuita en \texttt{claude.ai}. Bueno para análisis profundo, ensayos, código. Anthropic se enfoca en \emph{seguridad} de IA. \textbf{(3) Gemini (Google).} Integrado con servicios Google. Acceso en \texttt{gemini.google.com}. Bueno para tareas que se conectan con Gmail, Drive, Maps. Versión gratuita incluida.}
\milcpaso{2}{milcTurquesa}{\textbf{Las 5 cosas que SÍ pueden hacer los LLMs.} (1) \textbf{Explicar conceptos.} Pídele que te explique \emph{``derivadas''}, \emph{``fotosíntesis''}, \emph{``constitución colombiana''}: te dará explicación clara. (2) \textbf{Ayudar a estructurar texto.} Si tienes un ensayo en borrador, te puede sugerir estructura mejor. (3) \textbf{Traducir.} Inglés-español, español-francés, cualquier par. Mejor que Google Translate para frases complejas. (4) \textbf{Resumir.} Le pegas un texto largo y te lo resume en 5 puntos. (5) \textbf{Generar ideas.} Lluvia de ideas: nombres para un proyecto, temas para investigar, formas de empezar un texto.}
\milcpaso{3}{milcTurquesa}{\textbf{Las 5 cosas que NO pueden hacer.} (1) \textbf{No tienen información en tiempo real.} Si preguntas \emph{``¿qué pasó hoy en las noticias?''}, te dirá ``no tengo acceso a noticias en tiempo real''. Excepción: versiones con búsqueda web. (2) \textbf{Alucinan.} Inventan datos con seguridad. Si pides \emph{``cita una obra de Borges que hable de Cartago''}, puede inventar una cita falsa. (3) \textbf{No recuerdan entre conversaciones} (en versión gratuita). Cada conversación nueva arranca de cero. (4) \textbf{No tienen contexto personal.} No saben tu nombre, tu colegio, tus calificaciones. Cada vez que abres conversación, eres anónimo. (5) \textbf{No pueden hacer cálculos exactos largos sin error.} Si pides \emph{``cuánto es 6738 × 9281''}, puede dar respuesta incorrecta. Para mate exacta, usa calculadora.}
\milcpaso{4}{milcTurquesa}{\textbf{Cómo conversar bien con un LLM (preparación para S6).} \emph{(a) Saluda y da contexto:} \emph{``Soy estudiante de 7° grado en Colombia''}. Eso ayuda al LLM a adaptar tono. \emph{(b) Sé específico en lo que pides:} \emph{``Explícame el teorema de Pitágoras con un ejemplo del fútbol''} es mejor que \emph{``explícame Pitágoras''}. \emph{(c) Si no entiendes la respuesta, pide más:} \emph{``No entendí, ¿lo puedes explicar más simple?''}. La IA puede iterar. \emph{(d) Verifica lo que dice} (acuerdo 1 de la S1). Si te da un dato importante, verifica en otra fuente. \emph{(e) Aplica criterio} (acuerdo 4): si la respuesta no convence, pídele que la justifique o cámbiale el enfoque.}

\begin{drawbox}{milcTurquesa}{3 LLMs + 5 SÍ pueden + 5 NO pueden}
\textbf{3 LLMs principales:} \\[1mm]
\textbf{1.} ChatGPT (OpenAI, chat.openai.com). \\[1mm]
\textbf{2.} Claude (Anthropic, claude.ai). \\[1mm]
\textbf{3.} Gemini (Google, gemini.google.com). \\[1mm]
\textbf{Las 5 que SÍ pueden:} explicar · estructurar · traducir · resumir · generar ideas. \\[1mm]
\textbf{Las 5 que NO pueden:} info tiempo real · alucinan · sin memoria · sin contexto personal · cálculos exactos largos.
\end{drawbox}

\begin{softbox}{milcMostaza}{milcGris}{Errores comunes y su corrección}
\textbf{Errores frecuentes al usar LLMs por primera vez.} (1) \emph{Pedirle solo ``hola''} y esperar conversación profunda --- empieza con tarea concreta. (2) \emph{Aceptar la primera respuesta sin verificar} --- los LLMs alucinan a veces, sobre todo en datos específicos. (3) \emph{Darle datos personales sensibles} (cédula, dirección) --- recordatorio del acuerdo 2 de la S1. (4) \emph{Usarlo para hacer la tarea en lugar de aprender} --- te quita el aprendizaje, eso es responsabilidad ética (lo verás en S8 y S9). (5) \emph{Creer que lo que sabe ChatGPT es ``la verdad''} --- es ``la mejor predicción según sus datos'', no la verdad absoluta.
\end{softbox}

\begin{infoband}{milcTurquesa}


\textbf{Lo que va al cuaderno (Actividad 2 · EXPLICA).} Título: «Actividad 2 --- 3 LLMs + 5 sí pueden + 5 no pueden». \textbf{Formato:} bloque A con tabla 3x3 (LLM + URL + característica clave). Bloque B con lista de 5 capacidades. Bloque C con lista de 5 limitaciones. \textbf{Extensión:} 13 líneas. \textbf{Sabes que terminaste cuando} podrías recomendar qué LLM usar para qué tarea.
\end{infoband}

\puente{Con todo claro, abres un LLM gratuito y haces las 3 conversaciones reales.}

\milcestacion{milcMagenta}{Estación 3 de 3 · Hacer}

\begin{softbox}{milcMagenta}{milcGris}{Cómo vas a construir tu producto}
\textbf{Actividad 3 · APLICA --- 3 conversaciones reales con un LLM} (35 min · individual). Vas a abrir un LLM gratuito y conversar 3 veces con él. Sé honesto sobre tus dudas y aprende a evaluar las respuestas.
\end{softbox}

\milcpaso{1}{milcMagenta}{\textbf{Paso 1 --- Abre un LLM gratuito.} Decide cuál usar. Recomendaciones para 7° grado: \emph{Claude} (en \texttt{claude.ai}, no requiere cuenta para 5 mensajes), \emph{ChatGPT} (en \texttt{chat.openai.com}, requiere cuenta gratis), \emph{Gemini} (en \texttt{gemini.google.com}, con cuenta Google). Inicia con el que tengas más fácil acceso.}
\milcpaso{2}{milcMagenta}{\textbf{Paso 2 --- Conversación 1: Aprender algo difícil.} Toma el concepto que planeaste en Actividad 1. Escribe en el LLM: \emph{``Soy estudiante de 7° grado en Colombia. Necesito que me expliques [concepto] de manera sencilla, con un ejemplo concreto que pueda entender. Si después no entiendo, te pido que lo expliques distinto''}. Lee la respuesta. Si no entiendes, pide otra explicación. Conversa hasta que entiendas. \textbf{Anota en el cuaderno:} la pregunta inicial + 2-3 líneas de la respuesta clave + ¿entendiste mejor? ¿qué te ayudó?.}
\milcpaso{3}{milcMagenta}{\textbf{Paso 3 --- Conversación 2: Pedir recomendaciones.} Escribe: \emph{``Soy estudiante de 7° grado. Me gusta [género]. Recomiéndame 3 libros que me podrían gustar, con 1 línea de descripción cada uno''}. Recibe la respuesta. \textbf{Importante:} verifica al menos 1 de los libros recomendados. Búscalo en Google o en biblioteca. ¿Existe? Esto te entrena en \emph{verificación} (acuerdo 1).}
\milcpaso{4}{milcMagenta}{\textbf{Paso 4 --- Conversación 3: Redactar un correo formal.} Escribe: \emph{``Necesito redactar un correo formal al profe de Tecnología pidiéndole [tu pregunta]. ¿Me puedes ayudar con el borrador? Quiero que tenga las 5 partes formales''}. Recibe el borrador. \textbf{Importante:} no lo copies tal cual. Léelo. Adáptalo a tu voz. Cambia cosas que no suenan a ti. Esto practica el acuerdo 5 (voz propia).}
\milcpaso{5}{milcMagenta}{\textbf{Paso 5 --- Tabla comparativa + análisis.} Si tu colegio te lo permite, repite las 3 conversaciones con OTRO LLM distinto (Claude vs ChatGPT, por ejemplo). Compara las respuestas: ¿cuál fue más clara? ¿más profunda? ¿más útil? Anota en el cuaderno una tabla 3x2 (LLM x conversación) con tu evaluación. Si no puedes probar 2, evalúa solo el primero.}
\milcpaso{6}{milcMagenta}{\textbf{Paso 6 --- Reflexión personal + firma.} Al final escribe: \emph{``Mi primera vez con un LLM fue \_\_\_. Lo que más me sorprendió fue \_\_\_. Lo que verifiqué y resultó cierto: \_\_\_. Lo que NO me convenció: \_\_\_. Aplicaría las 5 acuerdos de la S1 así: \_\_\_''}. Firma con nombre y fecha.}

\begin{drawbox}{milcMagenta}{Antes de cerrar la S5}
\checkbox Tuve 3 conversaciones con un LLM gratuito. \\[1mm]
\checkbox Anoté las 3 conversaciones (preguntas + extractos de respuestas). \\[1mm]
\checkbox Verifiqué al menos 1 dato de las respuestas. \\[1mm]
\checkbox Apliqué los 5 acuerdos del uso responsable. \\[1mm]
\checkbox Hay reflexión personal de mi primera experiencia con LLM. \\[1mm]
\checkbox Está firmado con nombre y fecha.
\end{drawbox}

\begin{softbox}{milcMagenta}{milcGris}{Plantilla de trabajo}
\textbf{Estructura.} \textit{Paso 1:} abrir LLM. \textit{Paso 2:} conversación 1 (aprender concepto). \textit{Paso 3:} conversación 2 (recomendaciones + verificar). \textit{Paso 4:} conversación 3 (redactar + adaptar a voz propia). \textit{Paso 5:} tabla comparativa (opcional). \textit{Paso 6:} reflexión + firma.
\end{softbox}

\begin{infoband}{milcMagenta}


\textbf{Lo que va al cuaderno (Actividad 3 · APLICA).} Título: «Actividad 3 --- Mis 3 primeras conversaciones con un LLM». \textbf{Formato:} 3 bloques de conversación + tabla comparativa (opcional) + reflexión. \textbf{Extensión:} 1 página y media. \textbf{Sabes que terminaste cuando} tienes experiencia real con un LLM y opinión propia sobre su utilidad.
\end{infoband}

\puente{El producto es la tabla comparativa + 3 conversaciones capturadas + análisis.}

\milcestacion{milcMostaza}{Producto de la sesión}
\textbf{3 conversaciones reales con un LLM + verificación + reflexión · firmado}.
\begin{softbox}{milcMostaza}{milcGris}{Criterios de calidad}
(1) Tuve 3 conversaciones reales con un LLM gratuito. (2) Verifiqué al menos 1 dato de las respuestas. (3) Apliqué los 5 acuerdos del uso responsable. (4) El correo redactado fue adaptado a mi voz, no copiado directo. (5) Hay reflexión personal sobre mi primera experiencia.
\end{softbox}

\puente{Antes de salir, verifica que TÚ aplicaste los 5 acuerdos de la S1 (verificar, no datos personales, etc.).}

\milcestacion{milcVino}{Evaluación liberadora · cinco dimensiones}

\begin{softbox}{milcVino}{milcGris}{Sentido de la evaluación}
Evaluar no es solo revisar una nota. Es reconocer qué aprendí, qué cambió en mí, qué emociones manejé, cómo me relaciono con la ciudad y la comunidad, y qué le dejaré a quien viene detrás.
\end{softbox}

\textbf{Mi reflexión liberadora en cinco dimensiones.} Completa cada frase iniciada:

\small
\begin{tabularx}{\textwidth}{>{\bfseries\RaggedRight\arraybackslash}p{3.4cm}Y>{\RaggedRight\arraybackslash}p{4.6cm}}
\rowcolor{milcVino}\color{white}Dimensión & \color{white}\bfseries Mi respuesta (frame starter) & \color{white}\bfseries Algo que descubrí\\
1. Personal & Lo que más cambió en mí fue \linead{6.5cm} & \linead{4.2cm}\\
2. Emocional & La emoción más fuerte fue \linead{2.5cm} y la regulé \linead{3cm} & \linead{4.2cm}\\
3. Ciudadana & En democracia esto importa porque \linead{6cm} & \linead{4.2cm}\\
4. Local (Cartago) & En mi barrio sería útil para \linead{6.5cm} & \linead{4.2cm}\\
5. Intergeneracional & A mi abuela/abuelo le diría \linead{6.5cm} & \linead{4.2cm}\\
\end{tabularx}
\normalsize

\begin{infoband}{milcVino}
\textbf{Para la próxima sustentación.} Vuelve a esta tabla en seis meses. Verás que las respuestas cambian — no porque lo aprendido cambie, sino porque tú cambias. La evaluación liberadora es una conversación contigo a través del tiempo.
\end{infoband}


\puente{Cerramos con 3 voces sobre por qué dialogar con IA exige nuevas habilidades.}

\milcestacion{milcGradeDark}{Triángulo de pensamiento · cierre filosófico}

\begin{softbox}{milcVino}{milcGris}{Dussel · lente del nosotros}
\textit{``Dialogar con la IA no es someterse a ella. Es ejercer pensamiento crítico mientras conversas: preguntar, verificar, decidir tú.''}

Hay gente que abre ChatGPT y \textbf{copia la primera respuesta} sin pensar. Eso no es dialogar; es someterse. \emph{Dialogar críticamente} es: \textbf{leer con cuidado, hacer preguntas de seguimiento, verificar, contrastar con tu criterio, decidir si la respuesta es buena o no}. Las personas que dominan ese diálogo se vuelven más sabias. Las que copian sin pensar se vuelven dependientes. \emph{Hoy diste el primer paso en esa habilidad de toda la vida}.

\textbf{¿Voy a entrenar mi pensamiento crítico al usar IA, o a dejar que la IA piense por mí?}
\end{softbox}
\linead{\linewidth}\\[1mm]
\linead{\linewidth}

\begin{softbox}{milcMostaza}{milcGris}{Estoicismo · Marco Aurelio · cuidado interior}
\textit{``Cuando recibas una respuesta, antes de aceptarla, pregúntate: ¿es verdadera? ¿es útil para mí en mi situación? El sabio juzga antes de seguir.''}

Los LLMs dan respuestas con \emph{seguridad}: suenan correctas aunque a veces no lo sean. El estoico sabe que \textbf{seguridad no es verdad}. Antes de actuar según una respuesta, pasa por 2 filtros: \emph{(1) ¿es verdadera? (verifica con otra fuente)}, \emph{(2) ¿es útil para mí? (mi contexto puede ser distinto)}. Esos 2 filtros mentales en segundos te ahorran horas de errores.

\textbf{¿Estoy aplicando estos 2 filtros (verdad + utilidad propia) o acepto respuestas sin pasar por ellos?}
\end{softbox}
\linead{\linewidth}\\[1mm]
\linead{\linewidth}

\begin{softbox}{milcTurquesa}{milcGris}{Floridi · lente de la infoesfera}
\textit{``La calidad de tu conversación con la IA refleja la calidad de tu pensamiento. Quien sabe preguntar, obtiene respuestas valiosas. Quien no, obtiene basura.''}

\textbf{Los LLMs son espejo de tu pensamiento}. Si haces preguntas vagas, recibes respuestas vagas. Si haces preguntas profundas y específicas, recibes respuestas profundas y específicas. \emph{Las próximas sesiones (S6 y S7) te enseñarán a hacer preguntas mejor (prompting)}. Pero la habilidad de fondo es \emph{tu capacidad de formular bien la pregunta}, que viene de pensar bien sobre lo que necesitas.

\textbf{¿Mis preguntas son vagas o específicas? ¿Cómo mejora la respuesta cuando la pregunta es mejor?}
\end{softbox}
\linead{\linewidth}\\[1mm]
\linead{\linewidth}

\begin{infoband}{milcNegro}
\textbf{Nota para el docente.} Las tres formulaciones del triángulo son la síntesis del autor de la guía, no citas textuales de los pensadores. Si vas a citarlos en clase, remite a la obra original.
\end{infoband}

\milcestacion{milcVino}{Mi autoevaluación · marca una casilla por fila}

{\small
\setlength{\extrarowheight}{2.2pt}
\begin{tabularx}{\textwidth}{>{\RaggedRight\arraybackslash\bfseries}p{3.0cm}YYY}
\rowcolor{milcVino}\color{white}\bfseries Criterio & \color{white}\bfseries Lo logré & \color{white}\bfseries En proceso & \color{white}\bfseries Necesito apoyo\\[.6mm]
Comprensión del tema & \milcopcion{Explico las ideas con mis palabras.} & \milcopcion{Reconozco las ideas, falta precisión.} & \milcopcion{Necesito repasar lo básico.}\\[.8mm]
\rowcolor{milcGris}Pensamiento computacional & \milcopcion{Usé los pilares al armar mi producto.} & \milcopcion{Usé algunos pilares.} & \milcopcion{Necesito releer la estación 2.}\\[.8mm]
Aplicación y producto & \milcopcion{Mi producto cumple los criterios.} & \milcopcion{Está armado, con ajustes pendientes.} & \milcopcion{Aún está en construcción.}\\[.8mm]
\rowcolor{milcGris}Reflexión 5 dimensiones & \milcopcion{Respondí cada una con honestidad.} & \milcopcion{Respondí algunas con detalle.} & \milcopcion{Mis respuestas son muy generales.}\\[.8mm]
Triángulo de pensamiento & \milcopcion{Conecté las 3 voces con mi trabajo.} & \milcopcion{Conecté al menos una.} & \milcopcion{No las relaciono con mi proceso.}\\[.8mm]
\end{tabularx}}

\vspace{3mm}
\textbf{Hoy entendí que:}\\[1mm]
\linead{\linewidth}\\[1mm]
\linead{\linewidth}

\textbf{Mi compromiso para la próxima sesión es:}\\[1mm]
\linead{\linewidth}\\[1mm]
\linead{\linewidth}

\begin{softbox}{milcMostaza}{milcGris}{Cita de despedida · Marco Aurelio}
\textit{``El obstáculo en el camino se vuelve el camino. Lo que estorba al camino se vuelve camino.''}\\[1mm]
Cada error de hoy, bien observado, se convierte en pieza del camino que pisarás mañana.
\end{softbox}

\small
\begin{tabularx}{\textwidth}{>{\bfseries}p{3.6cm}Y}
\rowcolor{milcGradeDark!12}Nombre completo: & \linead{10cm}\\
Fecha: & \linead{4cm} \quad Curso/grupo: \linead{4cm}\\
Mi firma: & \linead{9cm}\\
\end{tabularx}
\normalsize

\begin{infoband}{milcGradeDark}
\textbf{Para la próxima sesión.} Esta semana, haz 2 conversaciones más con el LLM que más te gustó. Practica los acuerdos. La próxima clase aprendes \textbf{prompting básico}: cómo formular preguntas para obtener mejores respuestas. Es la habilidad que diferencia a quien usa IA con criterio de quien recibe basura.
\end{infoband}

\end{document}
